\usepackage{fontspec}
\usepackage{polyglossia}
\setdefaultlanguage[numerals=eastern]{arabic}
\setotherlanguage{english}
\newfontfamily\arabicfont[Script=Arabic]{Arial}
\setmainfont{Arial}

%==================================================
% [جديد] لوحة الألوان: أزرق سماوي مع رمادي
%==================================================
\usepackage{xcolor}
\definecolor{GovBlue}{HTML}{1E88C8}    % أزرق سماوي: الهوية الأساسية
\definecolor{GovDeep}{HTML}{14608F}    % [جديد] أزرق أغمق: للعناوين والتفاصيل
\definecolor{GovSlate}{HTML}{22303C}   % رمادي مزرق داكن: النصوص
\definecolor{GovSilver}{HTML}{8A96A3}  % رمادي: الخطوط الثانوية
\definecolor{GovGray}{HTML}{6B7785}    % رمادي: النصوص الثانوية
\definecolor{GovLight}{HTML}{E4F0F9}   % خلفية فاتحة: صفوف الجداول ومربع القسم
\definecolor{GovTrack}{HTML}{D3DEE8}   % مسار شريط التقدم في التذييل
\definecolor{GovTint}{HTML}{EEF6FC}    % [جديد] خلفية فاتحة جداً خلف الرأس والتذييل
% للتوافق مع الأسماء السابقة
\colorlet{GovNavy}{GovSlate}
\colorlet{GovTeal}{GovBlue}
\colorlet{GovGold}{GovSilver}
\colorlet{GovRed}{GovBlue}
\colorlet{GovBlack}{GovSlate}

%==================================================
% دمج عدة صفوف في خلية واحدة داخل الجداول
%==================================================
\usepackage{multirow}
%==================================================
% إدراج الصور والتحكم في حجمها وتدويرها
%==================================================
\usepackage{graphicx}

%==================================================
% دعم الخيار [H] لتثبيت الجداول والأشكال في مكانها
%==================================================
\usepackage{float}

%==================================================
% ضبط هوامش الصفحة وأبعاد الورقة
%==================================================
\usepackage{geometry}

\geometry{
    a4paper,
    landscape,
    left=1cm,
    right=1cm,
    bottom=1.9cm,   % [معدّل] كان 1cm؛ لرفع التذييل عن حافة الورقة
    top=0.5cm,
    includehead,
    headheight=35pt,
    headsep=0pt,
    footskip=31pt   % [معدّل] كان 25pt
}

%==================================================
% التحكم في كابشن الجداول وزيادة الارتفاع عن الجدول
%==================================================
\usepackage{caption}
\captionsetup[table]{
    font={Large,bf},
    labelfont={color=GovSlate},    % [جديد]
    textfont={color=GovSlate},   % [جديد]
    skip=5pt
}


%==================================================
% Header & Footer
%==================================================

\usepackage{fancyhdr}
\usepackage{lastpage}   % [جديد] لعرض: صفحة X من Y
\usepackage{adjustbox}  % [جديد] تصغير نص الرأس تلقائياً إن زاد عن المساحة المتاحة
\usepackage{tikz}       % [جديد] للمربعات المستديرة في الرأس والتذييل
\usepackage{eso-pic}    % [جديد] لرسم خلفية فاتحة بعرض الصفحة خلف الرأس والتذييل

\setlength{\headheight}{26pt}   % [معدّل] كانت 20pt لتتسع لمربع القسم
\setlength{\headsep}{12pt}
\setlength{\footskip}{31pt}   % [معدّل] كان 25pt
\renewcommand{\footruleskip}{2pt}   % [جديد] مسافة أصغر بين خط التذييل وشريطه

%==================================================
% Section - Subsection بدون ترقيم
%==================================================

\makeatletter

\newcommand{\CurrentSection}{}
\newcommand{\CurrentSubsection}{}

\renewcommand{\sectionmark}[1]{%
    \gdef\CurrentSection{#1}%
    \gdef\CurrentSubsection{}%
    \markright{\CurrentSection}%
}

\renewcommand{\subsectionmark}[1]{%
    \gdef\CurrentSubsection{#1}%
    \markright{\CurrentSection\ - \CurrentSubsection}%
}

\makeatother

%==================================================
% [جديد] محتوى الرأس والتذييل (خلفية بيضاء + خطوط ملوّنة)
%==================================================
% نص اختياري في وسط التذييل (فارغ افتراضياً)
% مثال: \renewcommand{\FooterNote}{تقرير العام الدراسي}
\newcommand{\FooterNote}{}
% نص الطرف الآخر من التذييل (غيّره أو اجعله فارغاً {})
\newcommand{\FooterRightNote}{الكتاب السنوي للاحصاءات التعليمية 2026/2027}
% شعار وزارة التربية في منتصف الرأس: ضع الصورة بجانب ملف الـ Rmd باسم moe_logo.png
% (إن لم توجد الصورة فلا يظهر شيء ولا يحدث خطأ). لتغيير الاسم أو الحجم:
% \renewcommand{\GovLogoFile}{اسم_آخر.png}   \setlength{\GovLogoHeight}{1cm}
\newcommand{\GovLogoFile}{image/moe_logo.png}
\newlength{\GovLogoHeight}
\setlength{\GovLogoHeight}{0.95cm}
% [قديم] بقي معرّفاً للتوافق فقط
\newcommand{\GovLogo}{}
% شريط تقدم الصفحات في التذييل (يُقسم الخط إلى لونين): \GovProgresstrue لتفعيله
\newif\ifGovProgress
\GovProgressfalse
\newdimen\GovProgDim
% خلفية فاتحة بعرض الصفحة خلف الرأس والتذييل: \GovTintfalse لإيقافها
\newif\ifGovTint
\GovTinttrue
\newlength{\GovHeadBandH}
\newlength{\GovFootBandH}
\AtBeginDocument{%
    \setlength{\GovHeadBandH}{\dimexpr1in+\voffset+\topmargin+\headheight\relax}%
    \setlength{\GovFootBandH}{\dimexpr\paperheight-1in-\voffset-\topmargin-\headheight-\headsep-\textheight-\footskip+25pt\relax}%
}
\AddToShipoutPictureBG{%
    \ifGovTint
        \AtPageUpperLeft{\raisebox{-\GovHeadBandH}[0pt][0pt]{\color{GovTint}\rule{\paperwidth}{\GovHeadBandH}}}%
        \AtPageLowerLeft{\color{GovTint}\rule{\paperwidth}{\GovFootBandH}}%
    \fi
}
% عدد الصفحات الكلي للعرض (بدون رابط/إطار إن كانت hyperref محمّلة)
\newcommand{\GovLastPage}{\ifdefined\hypersetup\pageref*{LastPage}\else\pageref{LastPage}\fi}
\newcounter{GovTotal}   % عدد الصفحات الكلي كرقم خام (يُحدَّث بعد الترجمة الأولى)
\makeatletter
\AtEndDocument{%
    \clearpage
    \immediate\write\@auxout{\string\setcounter{GovTotal}{\number\numexpr\value{page}-1\relax}}%
}

% الرأس - يسار: اسم القسم الحالي داخل مربع مستدير (يختفي إن كان فارغاً)
\newcommand{\GovHeadLeft}{%
    \protected@edef\GovTmpMark{\rightmark}%
    \ifx\GovTmpMark\@empty\else
        \adjustbox{max width=0.24\headwidth}{%
            \tikz[baseline=(GovHL.base)]\node[fill=white, draw=GovBlue!45, line width=0.6pt, text=GovDeep, rounded corners=3pt,
                inner xsep=8pt, inner ysep=4pt, font=\normalsize\bfseries](GovHL){\RL{\nouppercase{\rightmark}}};}%
    \fi
}
\makeatother
% الرأس - وسط: شعار وزارة التربية (صندوق بارتفاع صفري فلا يزيد ارتفاع الرأس)
\newcommand{\GovHeadCenter}{%
    \IfFileExists{\GovLogoFile}{%
        \raisebox{-3pt}[0pt][0pt]{\includegraphics[height=\GovLogoHeight]{\GovLogoFile}}%
    }{}%
}
% الرأس - يمين: الجهة في سطرين (المديرية ثم القسم والشعبة)
\newcommand{\GovHeadRight}{%
    \raisebox{0pt}[17pt][8pt]{\adjustbox{max width=0.42\headwidth}{\bfseries\color{GovSlate}%
    \shortstack{\fontsize{10.5}{11}\selectfont \RL{المديرية العامة للتربية في محافظة النجف الاشرف}\\[-1pt]
    \fontsize{8.5}{9}\selectfont\color{GovDeep}\RL{قسم التخطيط التربوي\textcolor{GovBlue}{\ \textbullet\ }شعبة الاحصاء}}}}%
}
% التذييل - يسار: صفحة X من Y داخل مربع أزرق مستدير
\newcommand{\GovFootLeft}{%
    \tikz[baseline=(GovFL.base)]\node[fill=GovBlue, text=white, rounded corners=4pt,
        inner xsep=10pt, inner ysep=4.5pt, font=\normalsize\bfseries](GovFL){\RL{صفحة\ \thepage\ من\ \GovLastPage}};%
}
% التذييل - وسط
\newcommand{\GovFootCenter}{%
    \normalsize\bfseries\color{GovGray}\RL{\FooterNote}%
}
% التذييل - يمين
\newcommand{\GovFootRight}{%
    \normalsize\bfseries\color{GovDeep}\RL{\FooterRightNote}%
}

%==================================================
% الرأس والتذييل
%==================================================

\pagestyle{fancy}
\fancyhf{}

% الرأس - يسار
\fancyhead[L]{\GovHeadLeft}

% الرأس - وسط: الشعار
\fancyhead[C]{\GovHeadCenter}

% الرأس - يمين
\fancyhead[R]{\GovHeadRight}

% التذييل
\fancyfoot[L]{\GovFootLeft}

\fancyfoot[C]{\GovFootCenter}

\fancyfoot[R]{\GovFootRight}

\input{models//footer.tex}

% خطوط الرأس والتذييل
\renewcommand{\headrulewidth}{6pt}
\renewcommand{\footrulewidth}{6pt}

% [جديد] خط الرأس بعرض الورقة كاملة دون قطع: أزرق سميك ثم رمادي رفيع
\renewcommand{\headrule}{%
    \hbox to\headwidth{\hss\vbox{\offinterlineskip
        \hbox to\paperwidth{\color{GovBlue}\vrule width \paperwidth height 3pt depth 0pt}%
        \vskip 1.5pt
        \hbox to\paperwidth{\color{GovSilver}\vrule width \paperwidth height 1.5pt depth 0pt}%
    }\hss}%
    \vskip-\headrulewidth
}

% [جديد] خط التذييل بعرض الورقة كاملة دون قطع: رمادي رفيع ثم أزرق سميك
% (إن فُعّل شريط التقدم يتحول الخط الأزرق إلى أزرق + مسار فاتح)
\renewcommand{\footrule}{%
    \vskip-\footruleskip\vskip-\footrulewidth
    \ifnum\value{GovTotal}>0
        \GovProgDim=\dimexpr\paperwidth*\value{page}/\value{GovTotal}\relax
        \ifdim\GovProgDim>\paperwidth \GovProgDim=\paperwidth\fi
    \else
        \GovProgDim=\paperwidth
    \fi
    \hbox to\headwidth{\hss\vbox{\offinterlineskip
        \hbox to\paperwidth{\color{GovSilver}\vrule width \paperwidth height 1.5pt depth 0pt}%
        \vskip 1.5pt
        \ifGovProgress
            \hbox to\paperwidth{\color{GovBlue}\vrule width \GovProgDim height 3pt depth 0pt%
            \color{GovTrack}\vrule width \dimexpr\paperwidth-\GovProgDim\relax height 3pt depth 0pt}%
        \else
            \hbox to\paperwidth{\color{GovBlue}\vrule width \paperwidth height 3pt depth 0pt}%
        \fi
    }\hss}%
    \vskip\footruleskip
}

%==================================================
% الصفحة الأولى (plain)
%==================================================

\fancypagestyle{plain}{%
    \fancyhf{}

    % الرأس - يسار
    \fancyhead[L]{\GovHeadLeft}

    % الرأس - وسط: الشعار
    \fancyhead[C]{\GovHeadCenter}

    % الرأس - يمين
    \fancyhead[R]{\GovHeadRight}

    % التذييل
    \fancyfoot[L]{\GovFootLeft}

    \fancyfoot[C]{\GovFootCenter}

    \fancyfoot[R]{\GovFootRight}

    \input{footer.tex}

    \renewcommand{\headrulewidth}{6pt}
    \renewcommand{\footrulewidth}{6pt}
}

%==================================================
% تطبيق التنسيق على أول صفحة أيضاً
%==================================================

\AtBeginDocument{%
    \pagestyle{fancy}
    \thispagestyle{fancy}
}

%==================================================
% تغيير حجم السكشن
%==================================================
\usepackage{titlesec}
\titleformat*{\section}{\Large\bfseries\color{GovDeep}}
\titleformat*{\subsection}{\Large\bfseries\color{GovBlue}}
\titleformat*{\subsubsection}{\Large\bfseries\color{GovGray}}

%==================================================
% ادراج ملفات بي دي اف
%==================================================
\usepackage{pdfpages}

%==================================================
% التحكم في ارتفاع الجداول
%==================================================
\renewcommand{\arraystretch}{1}

%==================================================
% الحزم الخاصة بالجداول
%==================================================
\usepackage{color}
\usepackage{rotating}
\usepackage{tabularray}
\definecolor{Silver}{rgb}{0.752,0.752,0.752}